
% ============================================================
% TP3 - PARTIE B : CRITERE ORIENTE TERMINAL
% Questions B1 a B4 - Yassine
% ============================================================

\chapter{Partie B : Offloading LTE/Wi-Fi selon la distance}
\label{chap:partie-b-yassine}

\section{Introduction}

Cette partie porte sur la mise en place d'un mécanisme de
\textit{Wi-Fi offloading} piloté par l'équipement utilisateur
(\textit{User Equipment}, UE). Le principe consiste à faire
passer le trafic de l'UE par le réseau LTE ou par le réseau
Wi-Fi en fonction de sa distance par rapport au point d'accès
(AP).

Nous procédons progressivement. La question B1 met en place
l'architecture LTE/Wi-Fi et la mobilité de l'UE. La question B2
ajoute un trafic applicatif entre l'UE et un serveur distant.
La question B3 introduit la décision de bascule fondée sur la
distance et mesure le débit reçu. Enfin, la question B4
analyse les interruptions autour des changements d'interface.

% ============================================================
\section{Question 1 : Mise en place du scénario LTE/Wi-Fi}
% ============================================================

\subsection{Objectif}

L'objectif est de construire un scénario NS-3 comprenant :
\begin{itemize}
    \item une station de base LTE (\textit{eNodeB}) ;
    \item un point d'accès Wi-Fi 802.11n ;
    \item un UE possédant une interface LTE et une interface Wi-Fi ;
    \item un modèle de mobilité permettant à l'UE de se déplacer
    en ligne droite à vitesse constante.
\end{itemize}

Cette première étape vérifie la création des équipements,
l'installation des piles réseau et l'attribution des adresses IP.
Elle prépare le scénario qui sera utilisé dans les questions
suivantes.

\subsection{Architecture et paramètres}

Le scénario contient un eNodeB placé en $(0,50,0)$ et un point
d'accès Wi-Fi placé en $(0,0,0)$. L'UE démarre en $(-100,0,0)$
et se déplace suivant l'axe $x$ avec une vitesse de $5$ m/s.
La durée de simulation est de $40$ secondes.

Le réseau Wi-Fi utilise la norme IEEE 802.11n et le SSID
\texttt{offload-ap}. Les équipements Wi-Fi reçoivent leurs
adresses dans le sous-réseau \texttt{10.0.1.0/24}. L'adresse
LTE de l'UE est attribuée par l'infrastructure EPC.

\begin{table}[H]
    \centering
    \caption{Paramètres du scénario de la question B1}
    \label{tab:b1-parametres}
    \begin{tabular}{ll}
        \toprule
        \textbf{Paramètre} & \textbf{Valeur} \\
        \midrule
        Nombre d'eNodeB & 1 \\
        Nombre de points d'accès Wi-Fi & 1 \\
        Nombre d'UE & 1 \\
        Norme Wi-Fi & IEEE 802.11n \\
        SSID & \texttt{offload-ap} \\
        Position de l'eNodeB & $(0,50,0)$ m \\
        Position de l'AP & $(0,0,0)$ m \\
        Position initiale de l'UE & $(-100,0,0)$ m \\
        Vitesse de l'UE & $5$ m/s suivant $+x$ \\
        Durée de simulation & $40$ s \\
        Sous-réseau Wi-Fi & \texttt{10.0.1.0/24} \\
        \bottomrule
    \end{tabular}
\end{table}

\subsection{Implémentation}

Le réseau LTE est créé à l'aide de \texttt{LteHelper} et
\texttt{PointToPointEpcHelper}. L'eNodeB et l'UE reçoivent
leurs équipements LTE avec les méthodes
\texttt{InstallEnbDevice()} et \texttt{InstallUeDevice()}.
L'UE est ensuite attaché à la station de base avec
\texttt{Attach()}.

Pour le Wi-Fi, nous utilisons \texttt{WifiHelper},
\texttt{WifiMacHelper} et \texttt{YansWifiPhyHelper}.
L'UE est configuré comme station Wi-Fi (\texttt{StaWifiMac})
et le point d'accès comme \texttt{ApWifiMac}. Les deux
interfaces appartiennent au même réseau IP.

La mobilité de l'UE est définie par
\texttt{ConstantVelocityMobilityModel}. La position initiale
et le vecteur vitesse sont fixés respectivement à
\texttt{Vector(-100,0,0)} et \texttt{Vector(5,0,0)}.
L'eNodeB et l'AP utilisent quant à eux un modèle de position
constante.

\capture{img/B/B1-configuration-mobilite.png}
{Extrait du code configurant les positions et la mobilité de l'UE.}
{fig:b1-mobilite}

\subsection{Résultats et vérification}

L'exécution de \texttt{offloading-q1.cc} affiche les adresses
IP des deux interfaces de l'UE ainsi que celles du point
d'accès. Les valeurs obtenues sont présentées dans le
tableau~\ref{tab:b1-adresses}.

\begin{table}[H]
    \centering
    \caption{Adresses IP affichées à l'exécution de B1}
    \label{tab:b1-adresses}
    \begin{tabular}{ll}
        \toprule
        \textbf{Équipement / interface} & \textbf{Adresse IP} \\
        \midrule
        UE -- LTE & \texttt{7.0.0.2} \\
        UE -- Wi-Fi & \texttt{10.0.1.1} \\
        Point d'accès -- Wi-Fi & \texttt{10.0.1.2} \\
        \bottomrule
    \end{tabular}
\end{table}

L'UE démarre bien à $(-100,0,0)$ et possède une vitesse
de $(5,0,0)$ m/s. La simulation se termine correctement.
Ces résultats valident la création des équipements et
l'adressage des interfaces. À ce stade, cette question
ne constitue pas encore une mesure de transfert applicatif
de bout en bout.

\capture{img/B/B1-execution.png}
{Sortie terminal de la simulation B1 : adresses IP, position et vitesse de l'UE.}
{fig:b1-execution}

\subsection{Conclusion de B1}

Le scénario de base LTE/Wi-Fi est opérationnel. L'UE dispose
de deux interfaces réseau et se déplace à vitesse constante
à proximité du point d'accès. Cette infrastructure constitue
la base des expériences de trafic et de bascule réalisées
dans les questions suivantes.

% ============================================================
\section{Question 2 : Génération d'un trafic entre l'UE et un serveur distant}
% ============================================================

\subsection{Objectif}

La question B2 consiste à ajouter un serveur distant et un
trafic applicatif continu entre ce serveur et l'UE. L'objectif
est de disposer d'un flux mesurable afin d'observer les
performances du réseau et de préparer l'étude de l'offloading.

\subsection{Topologie et adressage}

Nous conservons l'eNodeB, l'AP et l'UE de la question B1.
Un serveur distant (\textit{RemoteHost}) est ajouté et relié
au PGW de l'EPC par une liaison point-à-point. Une deuxième
liaison point-à-point relie le PGW au point d'accès Wi-Fi.
Cette dernière permet au trafic destiné au réseau Wi-Fi
d'atteindre l'AP.

Les sous-réseaux utilisés sont :
\begin{itemize}
    \item \texttt{1.0.0.0/8} pour la liaison PGW--serveur distant ;
    \item \texttt{10.0.1.0/24} pour le réseau Wi-Fi ;
    \item \texttt{10.0.2.0/24} pour la liaison PGW--AP.
\end{itemize}

Le routage statique est configuré pour permettre au serveur
d'atteindre le réseau de l'UE. Dans cette configuration
initiale, la route par défaut de l'UE est configurée sur
l'interface LTE. Le passage par Wi-Fi sera déclenché dans
la question B3.

\capture{img/B/B2-topologie-adressage.png}
{Topologie avec le serveur distant, le PGW et la liaison de retour vers l'AP.}
{fig:b2-topologie}

\subsection{Configuration du trafic UDP}

Un serveur UDP est installé sur le serveur distant au port
\texttt{9000}. Du côté de l'UE, un client UDP génère des
paquets de taille $1024$ octets avec un intervalle de
$10$ ms. Le trafic est donc un flux UDP à débit constant
(\textit{Constant Bit Rate}, CBR).

Le débit applicatif théorique est calculé par :

\[
D_{\mathrm{nominal}}
=
\frac{1024 \times 8}{0,01}
=
819\,200\ \mathrm{bit/s}
=
0,8192\ \mathrm{Mbit/s}.
\]

Le module FlowMonitor est utilisé pour collecter les
statistiques du flux, notamment le nombre de paquets
transmis et reçus, le débit et le délai moyen. Des traces
LTE, une capture PCAP Wi-Fi et un fichier XML FlowMonitor
sont également générés par le programme.

\capture{img/B/B2-code-trafic.png}
{Extrait du code installant le serveur UDP et le client CBR.}
{fig:b2-code-trafic}

\subsection{Résultats FlowMonitor}

Pour le flux applicatif identifié par le Flow ID 5, les
résultats affichés sont les suivants :

\begin{table}[H]
    \centering
    \caption{Résultats du flux applicatif de la question B2}
    \label{tab:b2-resultats}
    \begin{tabular}{lr}
        \toprule
        \textbf{Indicateur} & \textbf{Valeur mesurée} \\
        \midrule
        Flow ID & 5 \\
        Adresse source & \texttt{7.0.0.2} \\
        Adresse destination & \texttt{1.0.0.2} \\
        Paquets transmis & 3800 \\
        Paquets reçus & 3800 \\
        Paquets perdus & 0 \\
        Débit calculé par FlowMonitor & 0,841336 Mbit/s \\
        Délai moyen & 21,9307 ms \\
        \bottomrule
    \end{tabular}
\end{table}

Le nombre de paquets reçus est égal au nombre de paquets
transmis et aucune perte n'est reportée pour ce flux.
Le délai moyen observé est d'environ $21,93$ ms.

Le débit calculé par FlowMonitor peut être légèrement
différent du débit applicatif nominal, car il dépend
de la durée effective utilisée dans le calcul statistique
et des instants de transmission et de réception retenus.

Il est important de noter que, dans cette configuration
de B2, la route par défaut de l'UE est encore LTE. Les
résultats de cette question constituent donc la référence
avant l'activation du mécanisme de bascule vers le Wi-Fi.

\capture{img/B/B2-resultats-flowmonitor.png}
{Sortie terminal présentant les statistiques FlowMonitor du flux UDP de B2.}
{fig:b2-flowmonitor}

\subsection{Conclusion de B2}

Le trafic UDP entre l'UE et le serveur distant est généré
et observé avec succès. FlowMonitor fournit les indicateurs
nécessaires à l'analyse. Le flux de référence traverse
initialement le réseau LTE ; la question suivante introduit
le changement dynamique de route vers le Wi-Fi.

% ============================================================
\section{Question 3 : Offloading fondé sur la distance}
% ============================================================

\subsection{Objectif}

L'objectif de B3 est de mettre en place une décision
d'offloading côté UE. Lorsque l'UE entre dans la zone
de couverture du point d'accès Wi-Fi, le trafic est
redirigé vers Wi-Fi. Lorsqu'il quitte cette zone, le
trafic revient vers LTE.

Le critère de décision est la distance entre l'UE et
l'AP. Le rayon de couverture utilisé dans cette simulation
est fixé à $30$ m.

\subsection{Principe de décision}

À chaque vérification, le programme calcule la distance
$d(t)$ entre l'UE et le point d'accès. La règle de décision
est :

\[
\begin{cases}
\text{LTE} \rightarrow \text{Wi-Fi}, & d(t)\leq 30\ \text{m},\\
\text{Wi-Fi} \rightarrow \text{LTE}, & d(t)>30\ \text{m}.
\end{cases}
\]

Le contrôleur vérifie cette condition toutes les $0,5$ s.
Lorsqu'un changement est nécessaire, l'ancienne route
par défaut de l'UE est supprimée puis une nouvelle route
est installée vers la passerelle de l'interface sélectionnée.
Cette opération modifie le routage IP de l'UE sans recréer
les équipements réseau.

\capture{img/B/B3-code-controleur.png}
{Extrait du contrôleur d'offloading : calcul de distance et mise à jour de la route par défaut.}
{fig:b3-controleur}

\subsection{Évolution de la position et instants de bascule}

L'UE démarre à $x=-100$ m et avance à $5$ m/s. L'AP se
trouve à l'origine. La distance entre les deux équipements
diminue jusqu'au passage à proximité de l'AP, puis augmente.

Le seuil géométrique de $30$ m est franchi une première
fois à l'instant théorique :

\[
t_{\mathrm{entrée}}
=
\frac{100-30}{5}
=
14\ \text{s}.
\]

Il est franchi une deuxième fois lorsque l'UE s'éloigne :

\[
t_{\mathrm{sortie}}
=
\frac{100+30}{5}
=
26\ \text{s}.
\]

Comme le contrôleur ne vérifie la distance que toutes
les $0,5$ s, les changements de route sont observés
aux instants suivants :

\begin{itemize}
    \item LTE vers Wi-Fi à $t=14,5$ s, pour une distance
    mesurée de $27,5$ m ;
    \item Wi-Fi vers LTE à $t=26,5$ s, pour une distance
    mesurée de $32,5$ m.
\end{itemize}

Le léger décalage par rapport aux instants théoriques
est lié à la périodicité de vérification du contrôleur.

\capture{img/B/B3-transitions-terminal.png}
{Messages du terminal indiquant les deux transitions et les distances correspondantes.}
{fig:b3-transitions}

\subsection{Mesure du débit reçu}

Le programme mesure le volume de données reçu par le
serveur applicatif. Le débit est calculé toutes les
$100$ ms à partir de la différence entre le nombre
d'octets reçus pendant deux échantillonnages successifs :

\[
D(t)=
\frac{8\left(B(t)-B(t-\Delta t)\right)}
{\Delta t \times 10^6},
\qquad \Delta t=0,1\ \text{s},
\]

où $B(t)$ représente le nombre total d'octets reçus
jusqu'à l'instant $t$. Le résultat est exprimé en Mbit/s.

Les mesures sont enregistrées dans le fichier
\texttt{offloading-throughput.csv}, avec le temps,
le débit mesuré et l'interface sélectionnée.

\capture{img/B/B3-debit.png}
{Débit reçu en fonction du temps, avec les changements d'interface LTE/Wi-Fi.}
{fig:b3-debit}

\subsection{Résultats FlowMonitor}

FlowMonitor distingue les paquets émis depuis l'adresse
LTE de l'UE et ceux émis depuis son adresse Wi-Fi.
Les deux flux applicatifs principaux donnent les résultats
suivants :

\begin{table}[H]
    \centering
    \caption{Statistiques des flux LTE et Wi-Fi pour B3}
    \label{tab:b3-flowmonitor}
    \begin{tabular}{lrr}
        \toprule
        \textbf{Indicateur} & \textbf{Flux LTE} & \textbf{Flux Wi-Fi} \\
        \midrule
        Flow ID & 5 & 6 \\
        Adresse source & \texttt{7.0.0.2} & \texttt{10.0.1.1} \\
        Paquets transmis & 2600 & 1200 \\
        Paquets reçus & 2600 & 1200 \\
        Paquets perdus & 0 & 0 \\
        Débit & 0,575651 Mbit/s & 0,840884 Mbit/s \\
        Délai moyen & 21,9308 ms & 20,2019 ms \\
        \bottomrule
    \end{tabular}
\end{table}

Les deux flux confirment que le trafic applicatif a
emprunté successivement les deux interfaces. Le contrôleur
a sélectionné LTE avant l'entrée dans la zone Wi-Fi,
Wi-Fi pendant le passage dans cette zone, puis LTE
après la sortie.

Les échantillons du fichier CSV autour des changements
indiquent un débit reçu de $0,8192$ Mbit/s dans les
intervalles observés avant et après les bascules. Aucune
chute durable du débit n'est donc visible à la résolution
de mesure de $100$ ms. Cela ne permet toutefois pas
d'exclure un événement transitoire plus court que cette
fenêtre.

\subsection{Conclusion de B3}

Le mécanisme d'offloading basé sur la distance fonctionne.
Le trafic passe de LTE à Wi-Fi lorsque l'UE entre dans
la zone de couverture, puis revient à LTE lorsqu'il
s'en éloigne. Les mesures FlowMonitor distinguent les
deux chemins utilisés et le fichier CSV permet de suivre
le débit au cours du temps.

% ============================================================
\section{Question 4 : Mesure de l'interruption de service lors des bascules}
% ============================================================

\subsection{Objectif et méthode de mesure}

L'objectif de cette expérience est d'évaluer l'impact des changements
d'accès réseau LTE/Wi-Fi sur la continuité du service. Nous mesurons
la durée pendant laquelle le débit reçu devient inférieur à 50\,\%
du débit nominal de l'application.

Le trafic applicatif est généré avec les paramètres suivants :
\begin{itemize}
    \item Taille des paquets UDP : 1024 octets ;
    \item Intervalle entre deux paquets : 10 ms ;
    \item Débit nominal : 0,8192 Mbit/s ;
    \item Seuil minimal de débit : 0,4096 Mbit/s.
\end{itemize}

Le débit est évalué à l'aide d'une fenêtre glissante de 100 ms,
recalculée toutes les 10 ms. Pour chaque instant $t$, le débit est
calculé à partir des octets reçus pendant l'intervalle
$[t-0,1\,;\,t[$ :

\[
D(t)=\frac{8 \times B(t)}{0,1 \times 10^6}
\]

où $B(t)$ représente le nombre d'octets reçus dans la fenêtre
considérée.

Le seuil utilisé pour détecter une éventuelle interruption est :

\[
D_{\text{seuil}}=0,5 \times D_{\text{nominal}}
=0,4096\ \text{Mbit/s}
\]

Une fenêtre est considérée comme étant sous le seuil lorsque
son débit mesuré est inférieur à cette valeur.

\capture{img/B/B4-code-mesures.png}
{Extrait du code implémentant la fenêtre glissante de 100 ms avec un pas de 10 ms.}
{fig:b4-code-mesures}

\subsection{Analyse des transitions}

Deux transitions sont étudiées dans le scénario :
\begin{itemize}
    \item LTE vers Wi-Fi à $t=14,5$ s, lorsque le terminal entre
    dans la zone de couverture du point d'accès ;
    \item Wi-Fi vers LTE à $t=26,5$ s, lorsque le terminal quitte
    la zone de couverture Wi-Fi.
\end{itemize}

Pour chaque transition, l'analyse porte sur une période de deux
secondes centrée sur l'instant de bascule.

\begin{table}[H]
    \centering
    \caption{Résultats de la mesure de l'interruption de service}
    \label{tab:b4-interruption}
    \begin{tabular}{|l|c|c|}
        \hline
        \textbf{Indicateur} & \textbf{LTE $\rightarrow$ Wi-Fi}
        & \textbf{Wi-Fi $\rightarrow$ LTE} \\
        \hline
        Instant de transition & 14,5 s & 26,5 s \\
        \hline
        Période analysée & [13,5 ; 15,5] s & [25,5 ; 27,5] s \\
        \hline
        Durée mesurée sous le seuil & 0 s & 0 s \\
        \hline
        Paquets transmis & 200 & 200 \\
        \hline
        Paquets reçus & 200 & 200 \\
        \hline
        Paquets non reçus & 0 & 0 \\
        \hline
        Délai maximal observé & 26 ms & 21 ms \\
        \hline
    \end{tabular}
\end{table}

\capture{img/B/B4-resultats-terminal.png}
{Résultats de l'analyse des deux transitions LTE/Wi-Fi dans le terminal.}
{fig:b4-resultats-terminal}

\subsection{Analyse des résultats}

Les résultats montrent qu'aucune fenêtre glissante de 100 ms
n'a présenté un débit inférieur au seuil de 0,4096 Mbit/s
pendant les deux périodes de transition étudiées.

Pour la transition LTE vers Wi-Fi, les 200 paquets transmis
pendant la période d'analyse ont tous été reçus. Le délai maximal
observé est de 26 ms.

De même, lors de la transition Wi-Fi vers LTE, les 200 paquets
ont été reçus sans perte, avec un délai maximal de 21 ms.

La durée estimée sous le seuil est donc de 0 seconde pour les
deux transitions. Cela indique que le mécanisme de bascule
maintient un débit supérieur au seuil dans les fenêtres mesurées.

Il faut toutefois préciser que cette méthode repose sur des
fenêtres de 100 ms. Elle ne permet donc pas d'exclure une
dégradation très brève, d'une durée inférieure à la fenêtre
d'observation, qui pourrait être masquée par le calcul du débit
moyen.

\capture{img/B/B4-csv.png}
{Extrait du fichier CSV contenant les mesures du débit, les interfaces et les délais.}
{fig:b4-csv}

\subsection{Évolution du débit pendant l'offloading}

La figure suivante présente l'évolution du débit reçu pendant
la simulation. Elle fait apparaître le seuil de 50\,\% du débit
nominal ainsi que les deux instants de bascule entre LTE et Wi-Fi.

\capture{img/B/offloading-throughput.png}
{Évolution du débit reçu pendant les transitions LTE/Wi-Fi.}
{fig:b4-throughput}

\subsection{Comparaison avec un handover intra-LTE}

Dans un handover intra-LTE, le terminal passe d'un eNodeB à
un autre tout en restant dans le réseau LTE. Le changement
d'accès étudié dans cette expérience est différent : le terminal
bascule entre deux technologies d'accès, LTE et Wi-Fi.

Les deux mécanismes ont pour objectif de maintenir la
connectivité pendant le déplacement du terminal. Cependant,
les résultats de cette expérience ne constituent pas une
comparaison expérimentale directe avec un handover intra-LTE,
car aucun scénario de handover entre deux eNodeB n'a été
exécuté dans cette partie.

\subsection{Conclusion de B4}

Cette expérience montre que les deux transitions LTE/Wi-Fi
et Wi-Fi/LTE se sont déroulées sans perte de paquets dans
les périodes analysées. Aucune fenêtre glissante de 100 ms
n'a présenté un débit inférieur à 50\,\% du débit nominal.

Le mécanisme d'offloading implémenté permet ainsi de conserver
un débit mesuré supérieur au seuil défini pendant les deux
transitions étudiées, tout en maintenant la réception des
paquets applicatifs.

% ============================================================
\section{Bilan de la partie B1--B4}
% ============================================================

Les questions B1 à B4 ont permis de construire un scénario
combinant LTE et Wi-Fi, d'y générer un trafic UDP continu,
puis de mettre en place un mécanisme d'offloading piloté
par la distance entre l'UE et le point d'accès.

Les résultats montrent que l'UE utilise LTE avant d'entrer
dans la zone Wi-Fi, bascule vers Wi-Fi à $14,5$ s, puis
revient vers LTE à $26,5$ s. FlowMonitor confirme
l'utilisation des deux interfaces et aucune perte globale
n'est rapportée pour les flux applicatifs de cette
exécution. L'analyse par fenêtres de $100$ ms ne détecte
pas d'interruption selon le seuil défini.

Cette première stratégie constitue une base simple de
sélection d'interface. Ses limites, notamment la possibilité
de bascules répétées à proximité de la frontière de
couverture, motivent l'étude de l'hystérésis et des
paramètres de mobilité présentée dans les questions
suivantes.

% ============================================================
% FIN B1-B4
% ============================================================
